\section{Singularity envelopes at arbitrary mass}
\label{sec:arbitrary-mass-singularity-envelopes}
We now allow an arbitrary potential in a big class, without fixing the positive-mass reference $\phi$ used above. Let $\theta$ be a smooth closed real $(1,1)$-form with big class on the connected compact K\"ahler manifold $X$, and put $V=V_\theta$. Write
\[
 M(q)=\int_X\langle\theta_q^n\rangle,
 \qquad \Omega=\int_X\omega^n.
\]
The positive-mass characterization of singularity envelopes does not directly apply when $M(q)=0$. Mixing $q$ with $V$ gives positive mass, but that mass may tend to zero. The following normalized tail estimate permits passage to the limit. Attribution and source conventions are discussed in \cref{sec:history-energy}.

\begin{proposition}\label{prop:positive-mass-singularity-envelope-limit}
For $u\in\PSH(X,\theta)$ with $u\leq0$, put
\[
 w_\varepsilon=(1-\varepsilon)u+\varepsilon V,
 \qquad \phi_\varepsilon=P_\theta[w_\varepsilon],
 \qquad 0<\varepsilon<1.
\]
Then $\phi_\varepsilon\downarrow P_\theta[u]$ pointwise and in $L^1(X)$ as $\varepsilon\downarrow0$.
\end{proposition}
\begin{proof}
Put $\phi=P_\theta[u]$. Since $u\leq V$, the potentials $w_\varepsilon$ decrease to $u$. The envelope calculus gives $\phi\leq\phi_\varepsilon\leq V$ and makes $\phi_\varepsilon$ decrease with $\varepsilon$. Its limit $q$ is therefore $\theta$-psh and satisfies $\phi\leq q\leq V$. These integrable bounds also give $L^1$ convergence to $q$. We must identify $q$.

For $0\leq k\leq n$ define the finite non-pluripolar measures
\[
 \nu_k=\langle\theta_u^k\wedge\theta_V^{n-k}\rangle,
 \qquad a_k=\nu_k(X).
\]
Positive multilinearity, \cref{prop:nonpluripolar-local-calculus}, gives
\[
 \begin{aligned}
 \mu_{w_\varepsilon}
   =\sum_{k=0}^n\binom nk(1-\varepsilon)^k\varepsilon^{n-k}\nu_k,
 \\
 m_\varepsilon:=M(w_\varepsilon)
   =\sum_{k=0}^n\binom nk(1-\varepsilon)^k\varepsilon^{n-k}a_k.
 \end{aligned}
\]
Here $a_0=M(V)>0$ by \cref{prop:positive-mixed-mass-bound}, so $m_\varepsilon\geq\varepsilon^na_0>0$. Mass preservation and \cref{thm:positive-mass-model-characterization} show that $\phi_\varepsilon$ is a model potential of mass $m_\varepsilon$.

All differences in the following sets are taken on the common finite locus. The omitted pole sets have zero mass for these measures and for $\omega^n$. For $A>0$ set
\[
 r(A)=\max_{a_k>0}\frac{\nu_k(\{u<V-A\})}{a_k}.
\]
This maximum is over a nonempty finite set. Since each $\nu_k$ gives the poles zero mass, continuity from above yields $r(A)\downarrow0$. The identity $w_\varepsilon-V=(1-\varepsilon)(u-V)$ and the positive expansion give the uniform estimate
\begin{equation}\label{eq:normalized-singularity-tail}
 \frac{\mu_{w_\varepsilon}(\{w_\varepsilon<V-A\})}{m_\varepsilon}
                  \leq r(A),\qquad 0<\varepsilon<1.
\end{equation}

For a fixed $A>0$ put
\[
 p_{\varepsilon,A}=P_\theta(w_\varepsilon+A,V),
 \qquad p_A=P_\theta(u+A,V).
\]
We have $w_\varepsilon\leq p_{\varepsilon,A}\leq\phi_\varepsilon$, so mixed mass monotonicity implies
$M(p_{\varepsilon,A})=m_\varepsilon$ and $p_{\varepsilon,A}\in\mathcal E(X,\theta,\phi_\varepsilon)$. Decreasing rooftop continuity, \cref{prop:rooftop-algebra}, gives $p_{\varepsilon,A}\downarrow p_A$ as $\varepsilon\downarrow0$.

Apply \cref{thm:bounded-density-relative-prescribed-measure} with reference $\phi_\varepsilon$ and constant density $m_\varepsilon/\Omega$. It supplies $\chi_\varepsilon\in\mathcal E^1(X,\theta,\phi_\varepsilon)$ such that
\[
 \mu_{\chi_\varepsilon}=(m_\varepsilon/\Omega)\omega^n,
 \qquad \sup_X\chi_\varepsilon=0.
\]
By \cref{lem:model-normalized-supremum}, $\chi_\varepsilon\leq\phi_\varepsilon$. Normalized global compactness, \cref{prop:normalized-global-psh-compactness}, gives a constant $C_1$ independent of $\varepsilon$ with
$\int_X(-\chi_\varepsilon)\omega^n\leq C_1$.

\Needspace{7\baselineskip}
Fix $\delta>0$ and $B>1$. Put $\tau=\min\{1/2,\delta/(4B)\}$ and
\[
 \begin{aligned}
  G_{\varepsilon,B}&=\{\chi_\varepsilon>\phi_\varepsilon-B\},\\
  \psi_\varepsilon&=(1-\tau)\phi_\varepsilon
                       +\tau\chi_\varepsilon-\delta/2,\\
  E_{\varepsilon,A,\delta}&=\{\phi_\varepsilon>p_{\varepsilon,A}+\delta\}.
 \end{aligned}
\]
Since $\phi_\varepsilon\leq0$, the complement of $G_{\varepsilon,B}$ lies in $\{\chi_\varepsilon\leq-B\}$ and has volume at most $C_1/B$. Moreover,
\[
 \chi_\varepsilon-\delta/2\leq\psi_\varepsilon
                      \leq\phi_\varepsilon-\delta/2.
\]
Mass monotonicity thus places $\psi_\varepsilon$ in $\mathcal E(X,\theta,\phi_\varepsilon)$. On $G_{\varepsilon,B}$ we have $\psi_\varepsilon>\phi_\varepsilon-3\delta/4$, whence
$E_{\varepsilon,A,\delta}\cap G_{\varepsilon,B}\subset\{p_{\varepsilon,A}<\psi_\varepsilon\}$.
Positive multilinearity and relative full-mass comparison, \cref{prop:relative-mixed-comparison}, yield
\[
 \tau^n\frac{m_\varepsilon}{\Omega}
       \int_{E_{\varepsilon,A,\delta}\cap G_{\varepsilon,B}}\omega^n
 \leq\int_{\{p_{\varepsilon,A}<\psi_\varepsilon\}}\mu_{\psi_\varepsilon}
 \leq\int_{\{p_{\varepsilon,A}<\psi_\varepsilon\}}\mu_{p_{\varepsilon,A}}.
\]
The rooftop contact bound, \cref{cor:rooftop-contact-inequality}, is
\[
 \mu_{p_{\varepsilon,A}}
 \leq\mathbf1_{\{p_{\varepsilon,A}=w_\varepsilon+A\}}\mu_{w_\varepsilon}
       +\mathbf1_{\{p_{\varepsilon,A}=V\}}\mu_V.
\]
On the comparison set the second term vanishes because
$\psi_\varepsilon\leq V-\delta/2$. Contact with the first obstacle implies
$w_\varepsilon<V-A-\delta/2$. Dividing by the positive mass $m_\varepsilon$ and using \eqref{eq:normalized-singularity-tail} therefore gives
\begin{equation}\label{eq:ceiling-volume-recovery}
 \int_{E_{\varepsilon,A,\delta}}\omega^n
       \leq\frac{C_1}{B}+\Omega\tau^{-n}r(A).
\end{equation}
The bound is uniform in $\varepsilon$, even when $m_\varepsilon\to0$.

Take any sequence $\varepsilon_j\downarrow0$. Pointwise convergence on the common finite locus gives
\[
 \{q>p_A+\delta\}\subset
                 \liminf_j E_{\varepsilon_j,A,\delta}.
\]
Fatou's lemma transfers \eqref{eq:ceiling-volume-recovery} to this set. Since $p_A\leq\phi$, the same bound holds for $\{q>\phi+\delta\}$. First let $A\to\infty$ with $B,\delta$ fixed, and then let $B\to\infty$. The set has zero volume. For every $\delta>0$ this proves $q\leq\phi$ almost everywhere. Uniqueness of canonical psh representatives gives $q=\phi$ everywhere, as required.
\end{proof}
\index{singularity envelope!positive-mass perturbation}

\begin{definition}\label{def:singularity-ceiling}
For $u\in\PSH(X,\theta)$ let $\mathcal F_u$ consist of the potentials $v\in\PSH(X,\theta)$ satisfying $v\leq0$, $u\preceq v$ and
\[
 \int_X\langle\theta_v^k\wedge\theta_V^{n-k}\rangle
       =\int_X\langle\theta_u^k\wedge\theta_V^{n-k}\rangle,
                         \qquad 0\leq k\leq n.
\]
The \emph{singularity ceiling} is $C_\theta(u)=(\sup_{v\in\mathcal F_u}v)^*$.
\end{definition}
The family contains $u-\sup_Xu$ and is bounded above by zero. Thus its regularized supremum is $\theta$-psh. Both $C_\theta(u)$ and $P_\theta[u]$ are unchanged when a constant is added to $u$.
\index{singularity ceiling}

\begin{theorem}\label{thm:arbitrary-mass-ceiling-envelope}
For every $u\in\PSH(X,\theta)$ in a big class,
\[
 C_\theta(u)=P_\theta[u],
 \qquad P_\theta[P_\theta[u]]=P_\theta[u].
\]
These identities include potentials of zero Monge--Amp\`ere mass.
\end{theorem}
\begin{proof}
Subtract a constant so that $u\leq0$, and write $\phi=P_\theta[u]$. The mixed mass preservation in \cref{prop:p-envelope-calculus} shows that $\phi\in\mathcal F_u$. Hence $C_\theta(u)\geq\phi$.

Let $v\in\mathcal F_u$. For $0<\varepsilon<1$ put
\[
 w_\varepsilon=(1-\varepsilon)u+\varepsilon V,
 \qquad z_\varepsilon=(1-\varepsilon)v+\varepsilon V.
\]
The defining mixed mass equalities and positive multilinearity give
$M(w_\varepsilon)=M(z_\varepsilon)>0$. Also $w_\varepsilon\preceq z_\varepsilon$. Set $\zeta_\varepsilon=P_\theta[z_\varepsilon]$, a positive-mass model potential. Then $w_\varepsilon\preceq\zeta_\varepsilon$ and $M(w_\varepsilon)=M(\zeta_\varepsilon)$. The full-mass characterization, \cref{prop:relative-envelope-full-mass-characterization}, implies
\[
 P_\theta[w_\varepsilon]=\zeta_\varepsilon
                                      =P_\theta[z_\varepsilon].
\]
The preceding proposition, applied to $u$ and $v$, gives $P_\theta[u]=P_\theta[v]$. Since $v\leq0$, the envelope definition gives $v\leq P_\theta[v]=\phi$. Taking the supremum and its upper regularization proves $C_\theta(u)\leq\phi$.

Finally, $\phi\leq\phi_\varepsilon$ in the preceding proposition. Monotonicity and the positive-mass model property give
\[
 P_\theta[\phi]\leq P_\theta[\phi_\varepsilon]
                                    =\phi_\varepsilon.
\]
Letting $\varepsilon\downarrow0$ yields $P_\theta[\phi]\leq\phi$. The reverse inequality follows from $\phi\leq0$ and the envelope definition. This proves idempotence.
\end{proof}
\index{singularity envelope!idempotence}
